
\paragraph{Unified evidence reliability estimation.}
Current systems often estimate relevance, faithfulness, safety, and confidence with separate modules. Robust RAG needs a unified evidence reliability model that accounts for relevance, sufficiency, consistency, provenance, freshness, and adversarial risk.

\paragraph{Causal understanding of context influence.}
We need better tools for measuring how each passage changes the final answer. Passage attribution, counterfactual context removal, hidden-state probing, and answer-flip analysis can help distinguish useful evidence from harmful distraction.

\paragraph{Conflict-aware generation.}
Fact-level conflict resolution remains immature. Future systems should explicitly represent claims, sources, timestamps, and confidence, and generate answers that expose unresolved conflict rather than hiding it.

\paragraph{Adaptive control under cost constraints.}
Retrieval, reranking, verification, and human review all have costs. Robust RAG should optimise risk--cost tradeoffs: when is extra retrieval worth it, when is abstention better, and when should a system ask the user for clarification?

\paragraph{Adversarially realistic benchmarks.}
Many attacks assume simplified access patterns, while many defenses assume narrow threat models. Benchmarks should specify attacker capabilities, corpus mutability, retriever visibility, access control, and deployment constraints.

\paragraph{Multimodal, agentic, and long-context RAG.}
Robustness becomes harder when evidence includes images, tables, code, web pages, memory, tools, or long-context documents. Agentic RAG introduces persistent memory and tool calls, expanding both the control problem and the attack surface.

\paragraph{Versioned evidence maps.}
The literature is changing quickly, especially around RAG attacks and defenses. The practical implication is modest but important: survey artifacts should be versioned and machine-readable. Our supplementary CSV files provide a first snapshot of paper categories, source strata, threat-model fields, benchmark dimensions, assumptions, and known limitations; maintaining a living survey would require public versioning, update cadence, and re-coding procedures beyond the scope of this submission.
